\chapter{Linear Algebra}\label{linear-algebra}
\chapterstatus{complete}

\section{Introduction}\label{linear-algebra:section-introduction}

This chapter develops linear algebra over a field: linear maps and matrices, duality,
determinants, eigenvalues and canonical forms, bilinear forms, inner product spaces
and the spectral theorem, and complex structures on real vector spaces. The last topic
is the linear algebra underlying Hodge theory: the decomposition
$V_\CC = V^{1,0} \oplus V^{0,1}$ of a real vector space with a complex structure is the
pointwise version of the decomposition of forms on a complex manifold
(Chapter~\ref{complex-manifolds}) and of the Hodge decomposition
(Chapter~\ref{hodge-decomposition}).

Bases and dimension were introduced in Section~\ref{rings:section-rank}. References are
\cite{halmos-fdvs}, \cite{axler-linear} and \cite[Chapters XIII--XV]{lang-algebra}.

\noindent\textbf{Prerequisites.} Chapter~\ref{groups} (Groups),
Chapter~\ref{rings} (Rings and Modules) and Chapter~\ref{fields} (Fields and Galois
Theory).

Throughout, $k$ is a field and vector spaces are $k$-vector spaces, unless stated
otherwise.

\section{Linear maps and dimension}\label{linear-algebra:section-dimension}

\begin{proposition}\label{linear-algebra:proposition-subspaces}
Let $V$ be a vector space and $W \subseteq V$ a subspace.
\begin{enumerate}
\item $W$ has a \emph{complement}: a subspace $U$ with $V = W \oplus U$ (that is, $W + U = V$ and
$W \cap U = 0$). Then $U \cong V/W$ and $\dim V = \dim W + \dim V/W$.
\item $\dim W \leq \dim V$, and if $\dim V$ is finite and $\dim W = \dim V$ then $W = V$.
\item If $U, W \subseteq V$ are finite-dimensional, $\dim(U + W) + \dim(U \cap W) = \dim U + \dim W$.
\end{enumerate}
\end{proposition}

\begin{proof}
(1) Extend a basis $B_W$ of $W$ to a basis $B$ of $V$ (Theorem~\ref{rings:theorem-bases-exist}) and
let $U$ be spanned by $B \setminus B_W$. The projection $V \to V/W$ restricts to an isomorphism
$U \to V/W$, and $B = B_W \sqcup (B \setminus B_W)$ gives the formula (as cardinals).
(2) The first claim follows from (1). If $\dim W = \dim V < \infty$, a basis of $W$ has $\dim V$
elements and extends to a basis of $V$, which by Theorem~\ref{rings:theorem-dimension} has no
further elements.
(3) Extend a basis of $U \cap W$ to bases of $U$ and of $W$; the union of these bases is a basis of
$U + W$.
\end{proof}

\begin{theorem}[Rank--nullity]\label{linear-algebra:theorem-rank-nullity}
For a linear map $f \colon V \to W$, $V/\ker f \cong \im f$, so
$\dim V = \dim \ker f + \dim \im f$. If $\dim V = \dim W < \infty$, then $f$ is injective if and
only if it is surjective.
\end{theorem}

\begin{proof}
The isomorphism is the first isomorphism theorem
(Proposition~\ref{rings:proposition-module-basics}); combine with
Proposition~\ref{linear-algebra:proposition-subspaces}(1). If $\dim V = \dim W = n$, then $f$ is
injective iff $\dim \im f = n$ iff $\im f = W$ by
Proposition~\ref{linear-algebra:proposition-subspaces}(2).
\end{proof}

\begin{counterexample}\label{linear-algebra:counterexample-shift}
The last statement fails in infinite dimension. On $V = k^{(\NN)}$ the \emph{shift}
$e_n \mapsto e_{n+1}$ is injective but not surjective, and $e_0 \mapsto 0$, $e_{n+1} \mapsto e_n$ is
surjective but not injective.
\end{counterexample}

\section{Matrices}\label{linear-algebra:section-matrices}

\begin{definition}\label{linear-algebra:definition-matrix}
Let $V$ and $W$ have ordered bases $\mathcal{B} = (v_1, \ldots, v_n)$ and $\mathcal{C} = (w_1, \ldots, w_m)$.
The \emph{matrix} of a linear map $f \colon V \to W$ is the $m \times n$ matrix
$[f]_{\mathcal{C}, \mathcal{B}} = (a_{ij})$ with $f(v_j) = \sum_i a_{ij} w_i$: its $j$-th column contains the
coordinates of $f(v_j)$.
\end{definition}

\begin{proposition}\label{linear-algebra:proposition-matrices}
The map $f \mapsto [f]_{\mathcal{C}, \mathcal{B}}$ is an isomorphism of vector spaces
$\Hom_k(V, W) \to M_{m \times n}(k)$, so $\dim \Hom_k(V, W) = mn$. Composition corresponds to matrix
multiplication: $[g \circ f]_{\mathcal{D}, \mathcal{B}} = [g]_{\mathcal{D}, \mathcal{C}}[f]_{\mathcal{C}, \mathcal{B}}$. If
$\mathcal{B}'$, $\mathcal{C}'$ are other bases and $P = [\id]_{\mathcal{B}, \mathcal{B}'}$,
$Q = [\id]_{\mathcal{C}, \mathcal{C}'}$, then $[f]_{\mathcal{C}', \mathcal{B}'} = Q^{-1}[f]_{\mathcal{C}, \mathcal{B}}P$.
In particular the matrices of an endomorphism in different bases are \emph{similar}:
$A' = P^{-1}AP$.
\end{proposition}

\begin{proof}
A linear map is determined by the images of a basis, which can be arbitrary
(Proposition~\ref{rings:proposition-free-module}); this gives bijectivity, and linearity is clear.
If $f(v_j) = \sum_i a_{ij}w_i$ and $g(w_i) = \sum_l b_{li}u_l$, then
$g(f(v_j)) = \sum_l (\sum_i b_{li}a_{ij}) u_l$. The change of basis formula follows by writing
$f = \id_W \circ f \circ \id_V$ and using that $[\id]_{\mathcal{C}', \mathcal{C}} = Q^{-1}$.
\end{proof}

\begin{definition}\label{linear-algebra:definition-rank}
The \emph{rank} of a linear map is $\dim \im f$. The \emph{column rank} of a matrix
$A \in M_{m \times n}(k)$ is the dimension of the span of its columns in $k^m$, that is, the rank of
the map $k^n \to k^m$ it defines, and its \emph{row rank} is the column rank of the transpose
$A^T$.
\end{definition}

\section{Duality}\label{linear-algebra:section-duality}

\begin{definition}\label{linear-algebra:definition-dual}
The \emph{dual} of $V$ is $V^* = \Hom_k(V, k)$. For $f \colon V \to W$, the \emph{transpose} is
$f^* \colon W^* \to V^*$, $\varphi \mapsto \varphi \circ f$. For a subspace $W \subseteq V$, the \emph{annihilator}
is $W^\perp = \{\varphi \in V^* : \varphi|_W = 0\}$. If $(v_1, \ldots, v_n)$ is a basis of $V$, the \emph{dual
basis} $(v_1^*, \ldots, v_n^*)$ is given by $v_i^*(v_j) = \delta_{ij}$.
\end{definition}

\begin{proposition}\label{linear-algebra:proposition-duality}
Let $V$, $W$ be finite-dimensional.
\begin{enumerate}
\item The dual basis is a basis of $V^*$; so $\dim V^* = \dim V$.
\item The map $\mathrm{ev} \colon V \to V^{**}$ of Example~\ref{categories:example-double-dual} is an
isomorphism.
\item In dual bases, the matrix of $f^*$ is the transpose of the matrix of $f$.
\item $\dim W^\perp = \dim V - \dim W$ for $W \subseteq V$, and $(\im f)^\perp = \ker f^*$.
\item $\operatorname{rank} f^* = \operatorname{rank} f$. Hence row rank equals column rank for every matrix.
\end{enumerate}
\end{proposition}

\begin{proof}
(1) $\varphi = \sum_i \varphi(v_i)v_i^*$ for $\varphi \in V^*$, and the coefficients are determined by evaluating on
the $v_j$.
(2) $\mathrm{ev}$ is injective: if $v \neq 0$, extend $v$ to a basis and take the corresponding dual
basis element. Both sides have dimension $n$; apply
Theorem~\ref{linear-algebra:theorem-rank-nullity}.
(3) If $f(v_j) = \sum_i a_{ij}w_i$ then $f^*(w_i^*)(v_j) = a_{ij}$, so $f^*(w_i^*) = \sum_j a_{ij} v_j^*$.
(4) Extend a basis $v_1, \ldots, v_r$ of $W$ to a basis of $V$; then $W^\perp$ has basis
$v_{r+1}^*, \ldots, v_n^*$. And $\varphi \in \ker f^*$ iff $\varphi \circ f = 0$ iff $\varphi$ vanishes on $\im f$.
(5) By (4) and rank--nullity for $f^*$,
$\operatorname{rank} f^* = \dim W - \dim \ker f^* = \dim W - \dim(\im f)^\perp = \dim \im f$. Apply (3).
\end{proof}

\begin{counterexample}\label{linear-algebra:counterexample-double-dual}
In infinite dimension $\mathrm{ev} \colon V \to V^{**}$ is injective (by the same argument) but not
surjective. Let $k = \QQ$ and $V = \QQ^{(\NN)}$, which is countable. Then $V^* \cong \QQ^{\NN}$, which is
uncountable (it contains the sequences with values in $\{0, 1\}$). Choosing a basis of $V^*$ gives
an injection $V^* \to V^{**}$ (send the basis to its dual family), so $V^{**}$ is uncountable, and
cannot be the image of the countable set $V$.
\end{counterexample}

\section{Determinants}\label{linear-algebra:section-determinants}

In this section $R$ is a commutative ring. For $A \in M_n(R)$ we write $A = (a_1 | \cdots | a_n)$ with
columns $a_j \in R^n$.

\begin{definition}\label{linear-algebra:definition-alternating}
A map $D \colon (R^n)^n \to R$ is \emph{multilinear} if it is $R$-linear in each argument, and
\emph{alternating} if $D(a_1, \ldots, a_n) = 0$ whenever $a_i = a_j$ for some $i \neq j$. The
\emph{determinant} of $A \in M_n(R)$ is
\[
\det A = \sum_{\sigma \in S_n} \operatorname{sgn}(\sigma)\, a_{\sigma(1)1} a_{\sigma(2)2} \cdots a_{\sigma(n)n},
\]
where $\operatorname{sgn}$ is the sign of Proposition~\ref{groups:proposition-sign}.
\end{definition}

\begin{theorem}\label{linear-algebra:theorem-determinant}
The determinant, as a function of the columns, is multilinear and alternating with
$\det I = 1$. Every multilinear alternating $D \colon (R^n)^n \to R$ satisfies
$D(a_1, \ldots, a_n) = \det(a_1 | \cdots | a_n) \cdot D(e_1, \ldots, e_n)$.
\end{theorem}

\begin{proof}
Multilinearity is clear, and $\det I = 1$ since only $\sigma = \id$ contributes. If $a_i = a_j$ with
$i \neq j$, let $\tau = (i\,j)$. The map $\sigma \mapsto \sigma\tau$ is a bijection from even to odd permutations,
and the terms for $\sigma$ and $\sigma\tau$ have the same product of entries (as columns $i$ and $j$ agree)
and opposite signs; so the sum vanishes.

Let $D$ be multilinear and alternating. First, $D$ changes sign when two arguments are
exchanged: expand $0 = D(\ldots, a + b, \ldots, a + b, \ldots)$. Hence
$D(e_{\sigma(1)}, \ldots, e_{\sigma(n)}) = \operatorname{sgn}(\sigma) D(e_1, \ldots, e_n)$ for $\sigma \in S_n$, writing $\sigma$ as a
product of transpositions. Writing $a_j = \sum_i a_{ij}e_i$ and expanding by multilinearity,
$D(a_1, \ldots, a_n) = \sum_{\phi} a_{\phi(1)1} \cdots a_{\phi(n)n} D(e_{\phi(1)}, \ldots, e_{\phi(n)})$ over all maps
$\phi \colon \{1, \ldots, n\} \to \{1, \ldots, n\}$. The terms with $\phi$ not injective vanish, and the others give
$\det(a_1|\cdots|a_n) D(e_1, \ldots, e_n)$.
\end{proof}

\begin{proposition}\label{linear-algebra:proposition-determinant-properties}
For $A, B \in M_n(R)$:
\begin{enumerate}
\item $\det(AB) = \det A \det B$;
\item $\det A^T = \det A$, so $\det$ is also multilinear and alternating in the rows;
\item (Laplace expansion) for each $j$, $\det A = \sum_i (-1)^{i+j}a_{ij}\det A_{ij}$, where $A_{ij}$ is
obtained by deleting row $i$ and column $j$;
\item the \emph{adjugate} $\operatorname{adj}A$, with $(i,j)$ entry $(-1)^{i+j}\det A_{ji}$, satisfies
$A \operatorname{adj}A = \operatorname{adj}A \, A = \det(A) I$;
\item $A$ is invertible in $M_n(R)$ if and only if $\det A \in R^\times$.
\end{enumerate}
\end{proposition}

\begin{proof}
(1) The map $(b_1, \ldots, b_n) \mapsto \det(Ab_1 | \cdots | Ab_n)$ is multilinear and alternating, with value
$\det A$ on $(e_1, \ldots, e_n)$; apply Theorem~\ref{linear-algebra:theorem-determinant}.
(2) Reindex the sum by $\sigma \mapsto \sigma^{-1}$, using $\operatorname{sgn}(\sigma^{-1}) = \operatorname{sgn}(\sigma)$.
(3) Both sides are multilinear and alternating in the columns other than $j$, and linear in
column $j$. Writing $a_j = \sum_i a_{ij}e_i$, it suffices to treat $a_j = e_i$. Moving row $i$ to the top and
column $j$ to the front by $i - 1$ and $j - 1$ adjacent exchanges changes the determinant by
$(-1)^{i+j}$ and gives a matrix with first column $e_1$, whose determinant is $\det A_{ij}$ by the
Leibniz formula (only permutations fixing $1$ contribute).
(4) The $(j,j')$ entry of $\operatorname{adj}(A)A$ is $\sum_i (-1)^{i+j}\det A_{ij} a_{ij'}$, which by (3) is the
determinant of $A$ with column $j$ replaced by column $j'$: this is $\det A$ if $j = j'$ and $0$
otherwise. The other product is handled with (2).
(5) If $AB = I$ then $\det A\det B = 1$. Conversely $(\det A)^{-1}\operatorname{adj}A$ is an inverse by (4).
\end{proof}

\begin{definition}\label{linear-algebra:definition-determinant-endomorphism}
The determinant of an endomorphism $f$ of a finite-dimensional vector space (or of a finitely
generated free module over $R$) is the determinant of its matrix in any basis; it is well defined
since $\det(P^{-1}AP) = \det A$. The group $\GL(V)$ of invertible endomorphisms contains
$\SL(V) = \{f : \det f = 1\}$, and $\GL_n(R) = M_n(R)^\times$.
\end{definition}

\section{Eigenvalues}\label{linear-algebra:section-eigenvalues}

In this section $V$ is a vector space of finite dimension $n$ and $f \colon V \to V$ is linear.

\begin{definition}\label{linear-algebra:definition-eigen}
A scalar $\lambda \in k$ is an \emph{eigenvalue} of $f$ if $\ker(f - \lambda) \neq 0$; the nonzero elements of
the \emph{eigenspace} $V_\lambda = \ker(f - \lambda)$ are \emph{eigenvectors}. The \emph{characteristic
polynomial} is $\chi_f(x) = \det(xI - A) \in k[x]$ for the matrix $A$ of $f$ in any basis; it is monic of
degree $n$ and independent of the basis. The \emph{trace} of $f$ is $\operatorname{tr} f = \sum_i a_{ii}$. The map
$f$ is \emph{diagonalisable} if $V$ has a basis of eigenvectors.
\end{definition}

\begin{proposition}\label{linear-algebra:proposition-eigen-basics}
\begin{enumerate}
\item The eigenvalues of $f$ are the roots of $\chi_f$ in $k$.
\item $\chi_f(x) = x^n - (\operatorname{tr} f)x^{n-1} + \cdots + (-1)^n\det f$; in particular the trace is
independent of the basis.
\item Eigenvectors with pairwise distinct eigenvalues are linearly independent.
\end{enumerate}
\end{proposition}

\begin{proof}
(1) $\lambda$ is an eigenvalue iff $\lambda - f$ is not injective iff $\det(\lambda I - A) = 0$, by
Theorem~\ref{linear-algebra:theorem-rank-nullity} and
Proposition~\ref{linear-algebra:proposition-determinant-properties}(5).
(2) In the Leibniz expansion of $\det(xI - A)$ only $\sigma = \id$ contributes to the coefficients of $x^n$ and
$x^{n-1}$, since any other permutation moves at least two indices; evaluate at $x = 0$ for the constant
term.
(3) Suppose $\sum_{i=1}^m c_i v_i = 0$ with $m$ minimal and all $c_i \neq 0$. Applying $f - \lambda_m$ gives
$\sum_{i<m} c_i(\lambda_i - \lambda_m)v_i = 0$, a shorter relation with nonzero coefficients.
\end{proof}

\begin{theorem}[Cayley--Hamilton]\label{linear-algebra:theorem-cayley-hamilton}
$\chi_f(f) = 0$.
\end{theorem}

\begin{proof}
Let $A$ be the matrix of $f$. By Proposition~\ref{linear-algebra:proposition-determinant-properties}(4)
over $R = k[x]$, $(xI - A)\operatorname{adj}(xI - A) = \chi_f(x)I$. The entries of $\operatorname{adj}(xI - A)$ have degree
$< n$, so we can write $\operatorname{adj}(xI - A) = \sum_{i=0}^{n-1} B_i x^i$ with $B_i \in M_n(k)$; write
$\chi_f = \sum_{i=0}^n c_i x^i$. Comparing coefficients of $x^i$: $B_{i-1} - AB_i = c_i I$ for $0 \leq i \leq n$, with
$B_{-1} = B_n = 0$. Multiplying the $i$-th equation by $A^i$ on the left and summing, the left sides
telescope to $0$, so $\sum_i c_i A^i = 0$.
\end{proof}

\begin{definition}\label{linear-algebra:definition-minimal-polynomial}
The \emph{minimal polynomial} $m_f$ is the monic generator of the ideal
$\{p \in k[x] : p(f) = 0\}$, which is nonzero by Theorem~\ref{linear-algebra:theorem-cayley-hamilton}.
So $m_f \mid \chi_f$.
\end{definition}

\begin{lemma}[Kernel lemma]\label{linear-algebra:lemma-kernel}
If $p, q \in k[x]$ are coprime, then $\ker (pq)(f) = \ker p(f) \oplus \ker q(f)$.
\end{lemma}

\begin{proof}
Write $1 = ap + bq$. For $v \in \ker(pq)(f)$, $v = a(f)p(f)v + b(f)q(f)v$, where the first term lies in
$\ker q(f)$ (since $q(f)a(f)p(f)v = a(f)(pq)(f)v = 0$) and the second in $\ker p(f)$. If
$v \in \ker p(f) \cap \ker q(f)$ then $v = a(f)p(f)v + b(f)q(f)v = 0$.
\end{proof}

\begin{theorem}\label{linear-algebra:theorem-diagonalisable}
$f$ is diagonalisable if and only if $m_f$ is a product of distinct monic linear factors. In that case
$V = \bigoplus_\lambda V_\lambda$ over the eigenvalues. Commuting diagonalisable endomorphisms are
simultaneously diagonalisable.
\end{theorem}

\begin{proof}
If $V$ has a basis of eigenvectors with distinct eigenvalues $\lambda_1, \ldots, \lambda_r$, then
$\prod_i (f - \lambda_i) = 0$, so $m_f$ divides $\prod_i (x - \lambda_i)$. Conversely if
$m_f = \prod_i (x - \lambda_i)$ with distinct $\lambda_i$, the kernel lemma gives
$V = \ker m_f(f) = \bigoplus_i \ker(f - \lambda_i)$. If $g$ commutes with $f$, it preserves each $V_\lambda$, and its
restriction to $V_\lambda$ has minimal polynomial dividing $m_g$, hence is diagonalisable; choose
eigenbases of these restrictions.
\end{proof}

\begin{counterexample}\label{linear-algebra:counterexample-not-diagonalisable}
The rotation $\begin{pmatrix} 0 & -1 \\ 1 & 0 \end{pmatrix}$ has characteristic polynomial $x^2 + 1$ and no eigenvalue
over $\RR$, although it is diagonalisable over $\CC$. The matrix
$N = \begin{pmatrix} 0 & 1 \\ 0 & 0 \end{pmatrix}$ has $m_N = x^2$ and is not diagonalisable over any field: its only
eigenvalue is $0$ and $\ker N$ is one-dimensional.
\end{counterexample}

\section{Canonical forms}\label{linear-algebra:section-canonical-forms}

Let $f \colon V \to V$ be as in the previous section. By Example~\ref{rings:example-modules}, $V$ is a
$k[x]$-module $V_f$ with $x$ acting as $f$; it is finitely generated and torsion, being annihilated by
$\chi_f$.

\begin{definition}\label{linear-algebra:definition-companion}
The \emph{companion matrix} of a monic $g = x^d + c_{d-1}x^{d-1} + \cdots + c_0$ is the $d \times d$ matrix $C(g)$
with $1$'s below the diagonal, last column $(-c_0, \ldots, -c_{d-1})^T$ and zeros elsewhere; it is the
matrix of multiplication by $x$ on $k[x]/(g)$ in the basis $1, x, \ldots, x^{d-1}$. The \emph{Jordan block}
$J_e(\lambda)$ is the $e \times e$ matrix with $\lambda$ on the diagonal, $1$ directly above the diagonal and zeros
elsewhere.
\end{definition}

\begin{lemma}\label{linear-algebra:lemma-companion}
The minimal and characteristic polynomials of $C(g)$ both equal $g$. The minimal and characteristic
polynomials of $J_e(\lambda)$ both equal $(x - \lambda)^e$.
\end{lemma}

\begin{proof}
On $k[x]/(g)$ a polynomial $p$ acts as zero iff $p \cdot 1 = p$ lies in $(g)$, so the minimal polynomial is
$g$. The characteristic polynomial is monic of degree $d$ and divisible by $g$
(Definition~\ref{linear-algebra:definition-minimal-polynomial}), hence equal to $g$. The matrix
$J_e(\lambda)$ is the matrix of $x$ on $k[x]/((x - \lambda)^e)$ in the basis
$(x - \lambda)^{e-1}, \ldots, (x - \lambda), 1$, and the same argument applies.
\end{proof}

\begin{theorem}[Rational canonical form]\label{linear-algebra:theorem-rational-canonical-form}
There are unique monic polynomials $f_1 \mid f_2 \mid \cdots \mid f_s$ of positive degree such that $f$ has, in
a suitable basis, the block diagonal matrix $C(f_1) \oplus \cdots \oplus C(f_s)$. Moreover $m_f = f_s$ and
$\chi_f = f_1 \cdots f_s$. Two endomorphisms (or matrices) are similar if and only if they have the same
$f_i$, the \emph{invariant factors}.
\end{theorem}

\begin{proof}
By Theorem~\ref{rings:theorem-structure-pid}, $V_f \cong \bigoplus_i k[x]/(f_i)$ with $f_1 \mid \cdots \mid f_s$ monic of
positive degree (there is no free part since $V_f$ is torsion), uniquely determined. Bases of the
summands as in Definition~\ref{linear-algebra:definition-companion} give the block form. An
isomorphism $V_f \cong V_g$ of $k[x]$-modules is the same as a linear isomorphism $h$ with
$hf = gh$, so similarity corresponds to isomorphism of modules. The polynomial $f_s$ kills every
summand and is the minimal polynomial of the last one, so $m_f = f_s$; the characteristic
polynomial of a block matrix is the product of those of the blocks, which gives $\chi_f$ by
Lemma~\ref{linear-algebra:lemma-companion}.
\end{proof}

\begin{theorem}[Jordan normal form]\label{linear-algebra:theorem-jordan}
Suppose $\chi_f$ splits into linear factors over $k$ (for instance, $k$ algebraically closed). Then $f$
has, in a suitable basis, a block diagonal matrix $J_{e_1}(\lambda_1) \oplus \cdots \oplus J_{e_r}(\lambda_r)$, and the
multiset of pairs $(\lambda_i, e_i)$ is uniquely determined by $f$.
\end{theorem}

\begin{proof}
By Corollary~\ref{rings:corollary-primary-decomposition}, $V_f \cong \bigoplus_i k[x]/(p_i^{e_i})$ with $p_i$
monic irreducible, uniquely. Each $p_i$ divides $\chi_f$, hence is $x - \lambda_i$. Use the bases of
Lemma~\ref{linear-algebra:lemma-companion}.
\end{proof}

\begin{example}\label{linear-algebra:example-jordan}
A nilpotent endomorphism ($f^m = 0$ for some $m$) over any field has $\chi_f = x^n$ (its only possible
eigenvalue is $0$, and $m_f \mid x^m$), so it is similar to a direct sum of blocks $J_e(0)$; the sizes of the
blocks form a partition of $n$, and they classify nilpotent endomorphisms up to similarity. Such
nilpotent operators, and their Jordan decompositions, organise the limiting behaviour of
variations of Hodge structure (Chapter~\ref{vhs}).
\end{example}

\section{Bilinear forms}\label{linear-algebra:section-bilinear}

\begin{definition}\label{linear-algebra:definition-bilinear-form}
A \emph{bilinear form} on $V$ is a bilinear map $B \colon V \times V \to k$. It is \emph{symmetric} if
$B(v, w) = B(w, v)$, \emph{alternating} if $B(v, v) = 0$ for all $v$, and \emph{nondegenerate} if
$v \mapsto B(v, -)$ is an isomorphism $V \to V^*$ (for $V$ finite-dimensional: if $B(v, w) = 0$ for all $w$ implies
$v = 0$). Its \emph{Gram matrix} in a basis $(v_i)$ is $(B(v_i, v_j))$. For a subspace $W$,
$W^{\perp_B} = \{v : B(v, w) = 0 \text{ for all } w \in W\}$. A \emph{quadratic form} is $q(v) = B(v, v)$ for
a symmetric $B$.
\end{definition}

\begin{lemma}\label{linear-algebra:lemma-orthogonal-complement}
Let $B$ be symmetric or alternating on a finite-dimensional $V$, and $W \subseteq V$ a subspace on which $B$
is nondegenerate. Then $V = W \oplus W^{\perp_B}$. An alternating form $B$ satisfies $B(v, w) = -B(w, v)$.
\end{lemma}

\begin{proof}
$W \cap W^{\perp_B} = 0$ by nondegeneracy on $W$. Given $v \in V$, the functional $B(v, -)|_W$ equals
$B(w, -)|_W$ for a unique $w \in W$, and then $v - w \in W^{\perp_B}$ (for alternating forms use
$B(-, w) = -B(w, -)$). The last statement follows from expanding $B(v + w, v + w) = 0$.
\end{proof}

\begin{theorem}\label{linear-algebra:theorem-orthogonal-basis}
Let $\operatorname{char} k \neq 2$ and $B$ symmetric on a finite-dimensional $V$. Then $V$ has a basis
$(v_i)$ with $B(v_i, v_j) = 0$ for $i \neq j$.
\end{theorem}

\begin{proof}
By induction on $\dim V$. If $B = 0$ any basis works. Otherwise some $q(v) \neq 0$: if $B(v, w) \neq 0$ then
by the polarisation identity $2B(v, w) = q(v + w) - q(v) - q(w)$ one of $q(v + w)$, $q(v)$, $q(w)$ is
nonzero. Then $B$ is nondegenerate on $kv$, so $V = kv \oplus (kv)^{\perp_B}$ by
Lemma~\ref{linear-algebra:lemma-orthogonal-complement}; apply induction to $(kv)^{\perp_B}$.
\end{proof}

\begin{counterexample}\label{linear-algebra:counterexample-char-two}
The hypothesis $\operatorname{char} k \neq 2$ is needed. On $\mathbb{F}_2^2$ the form
$B((a, b), (c, d)) = ad + bc$ is symmetric and nondegenerate, but alternating
($B(v, v) = 2ab = 0$). A basis with $B(v_i, v_j) = 0$ for $i \neq j$ would make $B$ zero.
\end{counterexample}

\begin{theorem}[Sylvester's law of inertia]\label{linear-algebra:theorem-sylvester}
Let $B$ be a symmetric bilinear form on a real vector space $V$ of dimension $n$. There is a basis in which
the Gram matrix is diagonal with $p$ entries $1$, $q$ entries $-1$ and $n - p - q$ entries $0$, and the
numbers $p$ and $q$ do not depend on the basis. The pair $(p, q)$ is the \emph{signature} of $B$.
\end{theorem}

\begin{proof}
Take an orthogonal basis (Theorem~\ref{linear-algebra:theorem-orthogonal-basis}) and rescale each
$v_i$ with $q(v_i) \neq 0$ by $|q(v_i)|^{-1/2}$ (Proposition~\ref{number-systems:proposition-roots}).
For uniqueness, $p$ is the largest dimension of a subspace on which $q$ is positive definite: $q$ is
positive definite on the span $P$ of the $v_i$ with $q(v_i) = 1$, and if $U$ is a subspace on which $q$ is
positive definite, then $U$ meets the span $N$ of the other basis vectors, on which $q \leq 0$, only in $0$;
so $\dim U \leq n - \dim N = p$. Similarly for $q$.
\end{proof}

\begin{theorem}\label{linear-algebra:theorem-symplectic-basis}
Let $\omega$ be a nondegenerate alternating form on a finite-dimensional $V$. Then $\dim V = 2m$ is even, and $V$
has a \emph{symplectic basis} $e_1, \ldots, e_m, f_1, \ldots, f_m$ with $\omega(e_i, f_j) = \delta_{ij}$ and
$\omega(e_i, e_j) = \omega(f_i, f_j) = 0$.
\end{theorem}

\begin{proof}
If $V \neq 0$, choose $e_1 \neq 0$, and $f_1$ with $\omega(e_1, f_1) = 1$ by nondegeneracy. Then $e_1, f_1$ are
independent (as $\omega(e_1, e_1) = 0$), and $\omega$ is nondegenerate on their span $W$, whose Gram matrix is
$\begin{pmatrix} 0 & 1 \\ -1 & 0 \end{pmatrix}$. By Lemma~\ref{linear-algebra:lemma-orthogonal-complement},
$V = W \oplus W^{\perp_\omega}$, and $\omega$ is nondegenerate on $W^{\perp_\omega}$. Apply induction.
\end{proof}

\section{Inner product spaces and the spectral theorem}\label{linear-algebra:section-inner-product}

\begin{definition}\label{linear-algebra:definition-inner-product}
Let $\mathbb{K} = \RR$ or $\CC$. An \emph{inner product} on a $\mathbb{K}$-vector space $V$ is a map
$\langle -, - \rangle \colon V \times V \to \mathbb{K}$, linear in the first variable, with
$\langle w, v \rangle = \overline{\langle v, w \rangle}$ and $\langle v, v \rangle > 0$ for $v \neq 0$. (Over $\CC$ it is then
conjugate-linear in the second variable; such a form is also called a \emph{Hermitian inner
product}.) We write $\|v\| = \sqrt{\langle v, v \rangle}$. A family is \emph{orthonormal} if
$\langle v_i, v_j \rangle = \delta_{ij}$.
\end{definition}

\begin{proposition}\label{linear-algebra:proposition-cauchy-schwarz}
$|\langle v, w \rangle| \leq \|v\|\|w\|$, and $\|v + w\| \leq \|v\| + \|w\|$.
\end{proposition}

\begin{proof}
If $w = 0$ this is clear. Otherwise let $u = v - \frac{\langle v, w \rangle}{\|w\|^2}w$; then $\langle u, w \rangle = 0$ and
$0 \leq \|u\|^2 = \|v\|^2 - |\langle v, w \rangle|^2/\|w\|^2$. For the triangle inequality,
$\|v + w\|^2 = \|v\|^2 + 2\operatorname{Re}\langle v, w \rangle + \|w\|^2 \leq (\|v\| + \|w\|)^2$.
\end{proof}

\begin{theorem}[Gram--Schmidt]\label{linear-algebra:theorem-gram-schmidt}
If $v_1, v_2, \ldots$ is a finite or countable linearly independent family, there is an orthonormal family
$u_1, u_2, \ldots$ with $\operatorname{span}(u_1, \ldots, u_j) = \operatorname{span}(v_1, \ldots, v_j)$ for all $j$. Every
finite-dimensional inner product space has an orthonormal basis, and for every finite-dimensional
subspace $W$, $V = W \oplus W^\perp$, where $W^\perp = \{v : \langle v, w \rangle = 0 \text{ for all } w \in W\}$.
\end{theorem}

\begin{proof}
Define recursively $u_j' = v_j - \sum_{i<j}\langle v_j, u_i \rangle u_i$, which is nonzero since
$v_j \notin \operatorname{span}(u_1, \ldots, u_{j-1})$, and $u_j = u_j'/\|u_j'\|$. For the last claim, with an
orthonormal basis $(u_i)$ of $W$, write $v = \sum_i \langle v, u_i \rangle u_i + (v - \sum_i \langle v, u_i\rangle u_i)$;
the second term lies in $W^\perp$, and $W \cap W^\perp = 0$ by positivity.
\end{proof}

\begin{definition}\label{linear-algebra:definition-adjoint}
Let $V$ be a finite-dimensional inner product space and $f \colon V \to V$ linear. The \emph{adjoint} $f^\dagger$ is
the unique linear map with $\langle f(v), w \rangle = \langle v, f^\dagger(w) \rangle$ for all $v, w$. The map $f$ is
\emph{self-adjoint} if $f^\dagger = f$, \emph{normal} if $ff^\dagger = f^\dagger f$, and \emph{unitary} (or
\emph{orthogonal}, if $\mathbb{K} = \RR$) if $f^\dagger f = \id$. The unitary operators of $\CC^n$ form the
\emph{unitary group} $\mathrm{U}(n)$, and the orthogonal operators of $\RR^n$ the \emph{orthogonal group}
$\mathrm{O}(n)$.
\end{definition}

\begin{lemma}\label{linear-algebra:lemma-adjoint}
The adjoint exists and is unique. In an orthonormal basis its matrix is the conjugate transpose of the
matrix of $f$. If $W$ is $f$-invariant then $W^\perp$ is $f^\dagger$-invariant.
\end{lemma}

\begin{proof}
With an orthonormal basis $(u_i)$, the map $f^\dagger(w) = \sum_i \langle w, f(u_i) \rangle u_i$ works, and a linear map
is determined by the values $\langle v, f^\dagger(w) \rangle$. Its matrix entries are
$\langle f^\dagger(u_j), u_i \rangle = \overline{\langle f(u_i), u_j\rangle}$. If $f(W) \subseteq W$ and $v \in W^\perp$, then
$\langle f^\dagger(v), w \rangle = \langle v, f(w) \rangle = 0$ for $w \in W$.
\end{proof}

\begin{theorem}[Spectral theorem]\label{linear-algebra:theorem-spectral}
Let $V$ be a finite-dimensional inner product space.
\begin{enumerate}
\item If $\mathbb{K} = \CC$ and $f$ is normal, then $V$ has an orthonormal basis of eigenvectors of $f$.
\item If $f$ is self-adjoint (over $\RR$ or $\CC$), then all its eigenvalues are real and $V$ has an orthonormal
basis of eigenvectors of $f$.
\end{enumerate}
\end{theorem}

\begin{proof}
For self-adjoint $f$ and an eigenvector $v$ with eigenvalue $\lambda$,
$\lambda\|v\|^2 = \langle f v, v \rangle = \langle v, fv \rangle = \bar\lambda\|v\|^2$, so $\lambda \in \RR$.

(1) By induction on $\dim V$. By Theorem~\ref{fields:theorem-fta}, $\chi_f$ has a root, so $f$ has an eigenvector
$v$. Since $f$ and $f^\dagger$ commute, $f^\dagger$ preserves the eigenspace $V_\lambda$, and its restriction there has
an eigenvector $u$; then $u$ is an eigenvector of both $f$ and $f^\dagger$. Then $W = \CC u$ is invariant under
both, so $W^\perp$ is invariant under $f^\dagger$ and $(f^\dagger)^\dagger = f$ (Lemma~\ref{linear-algebra:lemma-adjoint}),
and the restriction of $f$ to $W^\perp$ is normal. Apply induction and normalise $u$.

(2) Over $\CC$ this follows from (1). Over $\RR$, choose an orthonormal basis and let $A$ be the real symmetric
matrix of $f$. As a complex matrix acting on $\CC^n$ with the standard inner product, $A$ is self-adjoint, so
it has an eigenvalue, which is real, and hence a real root $\lambda$ of $\chi_f$; so $f$ has an eigenvector $u \in V$.
The line $\RR u$ and its orthogonal complement are invariant under $f = f^\dagger$, and we conclude by induction.
\end{proof}

\begin{counterexample}\label{linear-algebra:counterexample-spectral}
Normality is necessary in (1): the matrix $N$ of
Counterexample~\ref{linear-algebra:counterexample-not-diagonalisable} is not diagonalisable at all. Over $\RR$,
normal is not enough: the rotation of that counterexample is orthogonal, hence normal, but has no real
eigenvector.
\end{counterexample}

\begin{corollary}\label{linear-algebra:corollary-positive-definite}
A real symmetric (or complex Hermitian) matrix $A$ defines an inner product $\langle v, w \rangle = w^* A v$ if and only
if all its eigenvalues are positive. Every inner product on $\RR^n$ (or $\CC^n$) is of this form.
\end{corollary}

\begin{proof}
Diagonalise $A = U D U^*$ with $U$ unitary (Theorem~\ref{linear-algebra:theorem-spectral}); then
$w^*Av = (U^*w)^* D (U^*v)$, which is positive definite iff all diagonal entries of $D$ are positive. The
matrix of an inner product in the standard basis is Hermitian.
\end{proof}

\begin{proposition}[Singular value decomposition]\label{linear-algebra:proposition-svd}
Let $f \colon V \to W$ be a linear map between finite-dimensional inner product spaces over $\mathbb{K} = \RR$ or $\CC$, of rank $r$. There are orthonormal bases $(v_1, \ldots, v_n)$ of $V$ and $(w_1, \ldots, w_m)$
of $W$ and real numbers $\sigma_1 \geq \cdots \geq \sigma_r > 0$ with $f(v_i) = \sigma_iw_i$ for $i \leq r$ and $f(v_i) = 0$ for $i > r$. The $\sigma_i^2$ are the nonzero eigenvalues of $f^\dagger f$, so the
\emph{singular values} $\sigma_i$ are determined by $f$; the largest one is the operator norm $\|f\| = \sup_{\|v\| = 1}\|f(v)\|$.
\end{proposition}

\begin{proof}
The endomorphism $f^\dagger f$ of $V$ is self-adjoint and $\langle f^\dagger fv, v\rangle = \|f(v)\|^2 \geq 0$, so by Theorem~\ref{linear-algebra:theorem-spectral} there is an orthonormal basis $(v_i)$ of $V$ with $f^\dagger fv_i = \lambda_iv_i$
and $\lambda_i \geq 0$; order it so that $\lambda_1 \geq \lambda_2 \geq \cdots$. Then $\langle f(v_i), f(v_j)\rangle = \langle f^\dagger fv_i, v_j\rangle = \lambda_i\delta_{ij}$, so the $f(v_i)$ are pairwise orthogonal, $f(v_i) = 0$ exactly when
$\lambda_i = 0$, and the number of $i$ with $\lambda_i > 0$ is the rank $r$. Put $\sigma_i = \sqrt{\lambda_i}$ and $w_i = f(v_i)/\sigma_i$ for $i \leq r$; these are orthonormal, and we extend them to an orthonormal basis of
$W$ (Theorem~\ref{linear-algebra:theorem-gram-schmidt}). For $v = \sum_ic_iv_i$ of norm one, $\|f(v)\|^2 = \sum_{i \leq r}\sigma_i^2|c_i|^2 \leq \sigma_1^2$, with equality at $v = v_1$.
\end{proof}

\begin{example}\label{linear-algebra:example-svd}
For $f(x, y) = (x + y, 0)$ on $\RR^2$, $f^\dagger f$ has matrix $\begin{pmatrix}1 & 1\\ 1 & 1\end{pmatrix}$ with eigenvalues $2$ and $0$ and eigenvectors $(1, 1)/\sqrt2$ and $(1, -1)/\sqrt2$. So $\sigma_1 = \sqrt2$, $w_1 = (1, 0)$,
and $\|f\| = \sqrt2$, although the eigenvalues of $f$ itself are $1$ and $0$: for non-normal maps the operator norm is not the largest absolute value of an eigenvalue.
\end{example}

\section{Complex structures}\label{linear-algebra:section-complex-structures}

\begin{definition}\label{linear-algebra:definition-complex-structure}
Let $V$ be a real vector space. Its \emph{complexification} is $V_\CC = V \otimes_\RR \CC$, a complex vector space
with the \emph{complex conjugation} $\overline{v \otimes z} = v \otimes \bar z$, which is $\RR$-linear and
conjugate-linear over $\CC$. A \emph{complex structure} on $V$ is an $\RR$-linear map $J \colon V \to V$ with
$J^2 = -\id$.
\end{definition}

\begin{proposition}\label{linear-algebra:proposition-complex-structure}
Let $J$ be a complex structure on a finite-dimensional real vector space $V$.
\begin{enumerate}
\item $V$ becomes a complex vector space with $(a + bi)v = av + bJv$; hence $\dim_\RR V = 2\dim_\CC V$ is even.
\item Extend $J$ $\CC$-linearly to $V_\CC$. Then $V_\CC = V^{1,0} \oplus V^{0,1}$, where
$V^{1,0} = \ker(J - i)$ and $V^{0,1} = \ker(J + i)$, and $\overline{V^{1,0}} = V^{0,1}$.
\item The map $V \to V^{1,0}$, $v \mapsto \frac12(v - iJv)$, is an isomorphism of complex vector spaces from
$(V, J)$ to $V^{1,0}$, and $v \mapsto \frac12(v + iJv)$ is a conjugate-linear isomorphism $V \to V^{0,1}$.
\end{enumerate}
\end{proposition}

\begin{proof}
(1) The axioms of a module over $\CC = \RR[x]/(x^2 + 1)$ follow from $J^2 = -1$
(Example~\ref{rings:example-modules} and Exercise~\ref{rings:exercise-gaussian-quotient}); a complex basis
$(v_j)$ gives the real basis $(v_j, Jv_j)$.
(2) The polynomial $x^2 + 1 = (x - i)(x + i)$ kills $J$ on $V_\CC$, and the kernel lemma
(Lemma~\ref{linear-algebra:lemma-kernel}) gives the decomposition. Conjugation commutes with $J$ (which is
real), so it exchanges the $\pm i$-eigenspaces.
(3) $J(v - iJv) = Jv + iv = i(v - iJv)$, so the image lies in $V^{1,0}$, and $J$ on $V$ corresponds to
multiplication by $i$. The map is injective because $v$ is recovered as twice the real part, and both sides
have complex dimension $\frac12\dim_\RR V$. The second map is the conjugate of the first.
\end{proof}

\begin{proposition}\label{linear-algebra:proposition-compatible-metric}
Let $J$ be a complex structure on a real inner product space $(V, g)$ with $g(Jv, Jw) = g(v, w)$. Put
$\omega(v, w) = g(Jv, w)$ and $h = g - i\omega$. Then $\omega$ is a nondegenerate alternating form with
$\omega(Jv, Jw) = \omega(v, w)$, and $h$ is a Hermitian inner product on the complex vector space $(V, J)$, whose
real part is $g$. Conversely, $g$ and $J$ determine $\omega$, and $\omega$ and $J$ determine
$g(v, w) = \omega(v, Jw)$.
\end{proposition}

\begin{proof}
From $J$-invariance, $g(Jv, w) = g(J^2 v, Jw) = -g(v, Jw)$. Hence
$\omega(w, v) = g(Jw, v) = -g(w, Jv) = -\omega(v, w)$, so $\omega$ is alternating;
$\omega(Jv, Jw) = g(J(Jv), Jw) = g(Jv, w) = \omega(v, w)$ by $J$-invariance of $g$; and $\omega$ is
nondegenerate because $\omega(v, Jv) = g(Jv, Jv) > 0$ for $v \neq 0$. For $h$:
$h(v, w) = g(v, w) + i g(v, Jw)$. Then $h(Jv, w) = g(Jv, w) + ig(Jv, Jw) = -g(v, Jw) + ig(v, w) = ih(v, w)$, so $h$ is
complex linear in the first variable; $h(w, v) = g(v, w) - ig(v, Jw) = \overline{h(v, w)}$; and
$h(v, v) = g(v, v) > 0$ for $v \neq 0$ since $g(v, Jv) = -g(Jv, v) = -g(v, Jv)$ vanishes. Finally
$\omega(v, Jw) = g(Jv, Jw) = g(v, w)$.
\end{proof}

\begin{example}\label{linear-algebra:example-standard-structure}
On $\RR^{2n}$ with coordinates $(x_1, y_1, \ldots, x_n, y_n)$, the standard complex structure is
$J(\partial_{x_j}) = \partial_{y_j}$, $J(\partial_{y_j}) = -\partial_{x_j}$ (writing $\partial_{x_j}$, $\partial_{y_j}$ for the standard basis vectors),
which identifies $(\RR^{2n}, J)$ with $\CC^n$ via $x + iy$. The standard inner product is $J$-invariant, the
associated form satisfies $\omega(\partial_{x_j}, \partial_{y_j}) = g(\partial_{y_j}, \partial_{y_j}) = 1$, so the $\partial_{x_j}$, $\partial_{y_j}$
form a symplectic basis, and $h$ is the standard Hermitian inner product of $\CC^n$. The vectors
$\partial_{z_j} = \frac12(\partial_{x_j} - i\partial_{y_j})$ span $V^{1,0}$. These formulas reappear as the local model of Kähler
manifolds (Chapter~\ref{kahler}).
\end{example}

\section{Exercises}\label{linear-algebra:section-exercises}

\begin{exercise}\label{linear-algebra:exercise-trace}
Show that $\operatorname{tr}(AB) = \operatorname{tr}(BA)$ for $A \in M_{m \times n}(k)$, $B \in M_{n \times m}(k)$, and deduce that similar
matrices have the same trace. Show that if $\chi_f = \prod_i (x - \lambda_i)$ then $\operatorname{tr} f = \sum \lambda_i$ and
$\det f = \prod \lambda_i$.
\end{exercise}

\begin{exercise}\label{linear-algebra:exercise-jordan-count}
Suppose $\chi_f$ splits. Show that the number of Jordan blocks $J_e(\lambda)$ with $e \geq j$ in the Jordan form of $f$
is $\dim\ker(f - \lambda)^j - \dim\ker(f - \lambda)^{j-1}$.
\end{exercise}

\begin{exercise}\label{linear-algebra:exercise-simultaneous-triangular}
Let $V$ be a finite-dimensional complex vector space and $\mathcal{F}$ a set of pairwise commuting endomorphisms.
Show that there is a basis in which all elements of $\mathcal{F}$ are upper triangular.
\end{exercise}

\begin{exercise}\label{linear-algebra:exercise-polar}
Show that every invertible complex matrix is uniquely $A = UP$ with $U$ unitary and $P$ Hermitian with positive
eigenvalues.
\end{exercise}

\begin{exercise}\label{linear-algebra:exercise-symplectic-determinant}
Let $\omega$ be a nondegenerate alternating form on $V$ of dimension $2m$ and $\mathrm{Sp}(V, \omega)$ the group of
$f \in \GL(V)$ with $\omega(fv, fw) = \omega(v, w)$. Show that $\det f = \pm 1$ for $f \in \mathrm{Sp}(V, \omega)$. (It is in fact
always $1$; this can be shown using the Pfaffian, Chapter~\ref{multilinear-algebra}.)
\end{exercise}

\begin{exercise}\label{linear-algebra:exercise-complex-structures-compatible}
Show that the set of complex structures $J$ on $\RR^{2}$ compatible with the standard inner product and with
$\omega(e_1, e_2) > 0$ has exactly one element.
\end{exercise}
